\section{Related work}
\label{sec:related_work}

This section locates the manuscript at the intersection of financial dependence, RMT, filtered networks and volatility modelling. It distinguishes established tools from the contribution of the present study: an integrated B3 workflow and a clearly retrospective assessment of whether full-sample structural descriptors are associated with volatility-model performance.

\subsection{Financial dependencies and econophysics}

Econophysics studies financial markets using concepts and methods from statistical physics, nonlinear dynamics, complex systems and network science \cite{smolyak_havlin_2022_econophysics,vogl_2022_financial_chaos}. A central observation is that financial returns exhibit stylized facts that are difficult to reconcile with independent Gaussian models. Heavy tails, volatility clustering, slow decay in absolute-return autocorrelations, intermittent bursts of activity and collective co-movement patterns are recurrent findings across asset classes, countries and sampling frequencies \cite{mantegna_stanley_1999_book,cont_2001_empirical}. These regularities motivate methods that focus not only on the marginal distribution of returns, but also on dependency, synchronization and network structure.

Financial dependency analysis lies at the intersection of econophysics and financial economics. In financial economics, dependence is central to portfolio risk, diversification, contagion, factor structure and systemic vulnerability. In econophysics, the same dependence is often represented through correlation matrices, eigenspectra, filtered graphs and hierarchical structures. Network-based studies of stock markets have used this representation to connect correlation structure with market organization and temporal evolution \cite{chi_2010_network,kumar_2012_correlation}. Raddant and Di~Matteo \cite{raddant_dimatteo_2023} emphasize that correlations, causality measures, volatility spillovers and networks offer additional views of financial systems. We follow this perspective: the goal is not merely to estimate pairwise correlations, but to decompose the correlation matrix, filter its noisy components, represent the remaining structure as networks and evaluate whether the extracted network structure has forecasting relevance.

Cross-correlation dynamics during crises have received attention because market stress often increases synchronization across assets. Zheng et~al. \cite{zheng_2012_changes} showed that changes in cross-correlation structure can serve as early indicators of systemic instability. Billio et~al. \cite{billio_2012_econometric} proposed econometric measures of connectedness and systemic risk in financial sectors. Junior and Franca \cite{junior_2012_correlation} documented that correlations among global stock indices tend to increase during crises, reducing the effective dimensionality of international markets. These findings motivate the rolling-correlation and crisis-synchronization analysis undertaken in this paper, where market-wide co-movement is treated as an object of study instead of a nuisance.

The relevance of these studies for the present manuscript is twofold. First, they show that financial dependence is dynamic and crisis-sensitive. Second, they indicate that a single unfiltered correlation matrix may conceal different layers of structure: a broad market component, group or sectoral organization and residual noise. This motivates the use of Random Matrix Theory and filtered networks as additional tools for separating and interpreting these layers.

\subsection{Random Matrix Theory in financial correlations}

RMT-based finance was established as an empirical programme for deciding which parts of a large correlation matrix are reproducible rather than as a new asset-pricing theory. Laloux et~al. \cite{laloux_1999_noise,laloux_2000_random} reported that most eigenvalues of equity correlation matrices are compatible with a random benchmark and showed that removing this component can improve portfolio construction. Plerou et~al. \cite{plerou_2002_random} added an interpretation of the outlying eigenvectors: the leading one is a broad market mode, whereas some subleading modes are associated with sectors or other collective groups. These findings made spectral structure an empirical object in its own right, but they do not imply that every eigenvalue outside the random band has a unique economic meaning.

The methodological literature subsequently focused on how aggressively the noisy component should be treated. Bun et~al. \cite{bun_2017_cleaning} review eigenvalue clipping, shrinkage and rotationally invariant estimators, and stress that the choice of cleaning rule affects both the stability of the estimated covariance matrix and its downstream use. The broader mathematical foundations are reviewed by Guhr et~al. \cite{guhr_muller_1998}. Thus, RMT studies differ in their objective: some seek a portfolio-optimal covariance estimator, while others use eigenvectors to identify collective market factors.

Our use is narrower and explicitly comparative. We use the Mar\v{c}enko--Pastur edge to organize the empirical spectrum, remove the dominant market mode, retain candidate group modes and preserve the remaining trace in a residual component. The resulting matrix is not presented as an optimal covariance estimator. It is one controlled input to the paired MST/PMFG comparison, allowing us to ask whether localized topology and volatility-model performance change after the common market component is removed. Because the decomposition is estimated from the complete synchronized panel, the comparison is retrospective.

\subsection{Financial networks: MST, PMFG and filtered graphs}

Mantegna's introduction of the MST made it possible to study market organization without imposing a fixed correlation threshold \cite{mantegna_1999_hierarchical}. Subsequent applications used the tree to describe sectoral and geographical taxonomies, portfolio organization and time-varying market integration. Onnela et~al. \cite{onnela_2003_dynamics} documented the evolution of market correlations and the changing concentration of the tree, while Bonanno et~al. \cite{bonanno_2004_networks} and Song et~al. \cite{song_2012_evolution} showed that the resulting hierarchies depend on both sector and market regime. The common result is interpretability, accompanied by a loss of cycles and alternative routes.

The PMFG line of work addresses that loss of local structure. Tumminello et~al. \cite{tumminello_2005_tool,tumminello_2010_correlation} used planar filtering to retain stronger links together with triangular relationships, making clustering and clique structure available for financial interpretation. Massara et~al. \cite{massara_2016_tmfg} proposed the TMFG as a related scalable construction for larger systems, while temporal-network formulations retain information about the evolution of market links \cite{zhao_2018_stock}. Other studies connected filtered topology to economic questions: Pozzi et~al. \cite{pozzi_2013_spread} related peripheral positions to diversification, Aste et~al. \cite{aste_2010_correlation} examined network changes in volatile periods, and Kenett et~al. \cite{kenett_2012_dominating} used partial correlations to distinguish direct from market-mediated links.

The gap relevant here is not the absence of filtering methods, but the limited integration of spectral and topological comparisons. Most network studies construct one graph from the observed correlations and interpret its hubs or communities. We instead construct matched MST and PMFG representations before and after market-mode removal, hold the asset universe and ranking procedure fixed, and compare centrality, clustering and sector retention. This design treats topology as a diagnostic of spectral filtering and as a retrospective descriptor for the forecasting experiment, rather than as evidence of causal links.

\subsection{Financial networks in emerging markets}

Evidence from emerging markets shows why results from large developed-market datasets cannot simply be transferred to B3. Institutional structure, sector concentration, liquidity and exposure to commodities or exchange rates can alter both the level and the persistence of cross-sectional dependence. Existing work also uses several different network objects. Tabak et~al. \cite{tabak_2010_topology} studied the Brazilian interbank network and its vulnerability to targeted failures, while Silva et~al. \cite{silva_2016_structure} examined the position and dynamics of emerging economies in the global financial network. Wang et~al. \cite{wang_2018_correlation} compared Pearson and partial-correlation networks across world markets and found that the inferred structure depends on the dependence measure.

Studies closer to the equity-market setting remain fragmented. Pereira et~al. \cite{pereira_2019_multiscale} used detrended cross-correlation networks for a set of stock markets that includes Brazil; Ferreira et~al. \cite{ferreira_2019_detrended} compared detrended dependencies between oil and stock markets; Leonel Caetano and Yoneyama \cite{leonel_2022_brazilian} analysed the Brazilian stock market with complex-network measures; and Stosi{\'c} et~al. \cite{stosic_2015_correlations} focused on Brazilian agrarian commodities and stocks. These contributions establish that Brazilian data contain nontrivial multiscale and network structure, but they do not jointly evaluate a broad B3 equity universe, RMT group-mode filtering, planar graphs and volatility forecasting.

Our contribution is therefore contextual rather than a claim that B3 is representative of all emerging markets. We use the original and RMT-filtered networks to separate broad market integration from more localized organization, and we report the comparison descriptively. The subsequent volatility exercise asks whether these ex-post structural differences are associated with model performance; it does not infer a universal emerging-market mechanism or a tradable signal.

\subsection{Volatility forecasting and machine learning}

The forecasting literature provides two relevant benchmarks. Andersen et~al. \cite{andersen_2003_modeling} established the empirical role of realized volatility, while Poon and Granger \cite{poon_granger_2003} reviewed the comparative performance of volatility models. The HAR formulation of Corsi \cite{corsi_2009_simple} is particularly relevant because it gives a transparent multi-horizon benchmark against which additional information can be judged. Forecast comparisons should be evaluated against strong conventional benchmarks rather than only against one another \cite{hansen_lunde_2005}. These studies motivate our use of conventional return and volatility predictors as the baseline rather than treating a flexible learner as the default.

Machine-learning work has expanded the feature space and the class of admissible nonlinearities. Bucci \cite{bucci_2020_realized} used neural networks for realized-volatility forecasting, Christensen et~al. \cite{christensen_2022_realized} developed a broader machine-learning framework, and Zhang and Broadstock \cite{zhang_2023_network_volatility} connected financial-network topology with market-risk modelling. Random forests, boosting and regularized linear models remain useful comparison points because they expose different bias--variance and interpretability trade-offs \cite{breiman_2001_random,friedman_2001_greedy,hoerl_kennard_1970_ridge}.

Two methodological issues are central in this literature. First, realized volatility is a noisy proxy for latent variance, so Patton \cite{patton_2011_volatility} recommends loss functions such as QLIKE that are robust to measurement error. Second, flexible models can produce spurious improvements when temporal ordering, feature selection or the information set is not respected. Our experiment addresses these issues with a chronological train--validation--test split, complementary losses and nested feature sets. The network variables are treated as an incremental structural augmentation: because Sets~B and~C use full-sample RMT and network descriptors, their results are interpreted as retrospective associations, not point-in-time forecast gains.
